\chapter{Aufstieg, Methodik, ELO-Verlauf}\label{kap:weitere}

\section{Die Aufstiegsseite der Regionalligen}

Die Seite \code{rl-aufstieg.html} gehört zu keiner Liga. Sie fasst den
Aufstieg aus allen fünf Staffeln zusammen
(\code{.render\_aufstiegsseite}, \datei{RCode/generate_static_site.R:860--904}).
Unter der Überschrift \enquote{Aufstieg in die 3.~Liga} erklärt ein kurzer
Text die Regel. Danach folgt je Staffel eine Tabelle mit drei Spalten:

\begin{description}[font=\sffamily\bfseries\small]
  \item[Meister] $P(\text{Platz 1})$ aus dem Hauptlauf der Staffel.
  \item[Aufstieg] Bei Direktaufsteigern gleich dem Meister, bei den
        Playoff-Staffeln nach Gleichung~\eqref{eq:rl-aufstieg}.
  \item[Siegquote] $P(\text{Aufstieg})/P(\text{Meister})$, nur bei den
        Playoff-Staffeln. Leer, wenn der Quotient nicht definiert ist.
\end{description}

Welche Staffel Playoff-Staffel ist, liest die Seite aus der Ansicht ab und
nicht aus festen Staffelnamen. So folgt sie der jährlichen Rotation ohne
Codeänderung. Gezeigt werden nur Mannschaften mit einer Meister- oder
Aufstiegschance über null.

\screenshot{aufstiegsseite}{Die Aufstiegsseite mit den fünf Staffeln.}

\section{Die Methodik-Seite}

Der Text der Methodik-Seite ist redaktionell und liegt als fertiges HTML in
\datei{RCode/site_assets/methodik_content.html}. Er geht auf den Blogbeitrag
\enquote{Was die Prognosen mit Schach zu tun haben} vom 1.~April 2015 zurück
und wurde für den Relaunch überarbeitet. Er erklärt das Modell ohne Formeln:

\begin{enumerate}
  \item Die Stärkeschätzung ELO, Heimvorteil 40 Punkte, gedämpfte Wirkung
        hoher Siege,
  \item vom Stärkeunterschied zum Ergebnis: Torerwartung, Poisson-Verteilung,
        das Beispiel Bayern gegen Stuttgart,
  \item zehntausend Saisons: das Monte-Carlo-Verfahren,
  \item Abstieg aus den Regionalligen und warum er erst am Saisonende
        feststeht,
  \item was die Prognose kann und was nicht: Häkchen, Gleichstand,
        Zweitvertretungen, Wertungen, Modellgüte.
\end{enumerate}

Ein farbig abgesetzter Absatz vermerkt die Neukalibrierung des
Heimvorteils von 65 auf 40 Punkte im September 2026.

\section{Die ELO-Verlaufsseiten}\label{sec:verlauf}

Seit Oktober 2026 hat jede Liga eine Verlaufsseite (Issue~\#184). Sie zeigt,
wie sich die Stärkeschätzung jeder Mannschaft über die Saison entwickelt hat.
Gestaltet ist sie nach den Grundsätzen von Edward Tufte, mit den
Gestaltungsmitteln der Website.

\screenshot{verlauf-bundesliga}{ELO-Verlauf mit hervorgehobenem Verein und Tooltip.}

\subsection*{Aufbau}

\begin{enumerate}
  \item Kopf und Navigation (die Liga-Links führen zu den anderen
        Verlaufsseiten), darunter ein Rücklink \enquote{← Prognose \emph{Liga}}.
  \item Überschrift \enquote{ELO-Verlauf \emph{Liga} \emph{Saison}} und der
        Einleitungssatz \enquote{Jede Linie ist ein Verein, jeder Knick ein
        Spiel. […] Die Zahl am Ende jeder Linie ist der heutige Stand.}
  \item Das Diagramm mit einem Umschalter \enquote{Nach Spielen} /
        \enquote{Nach Datum}.
  \item Ein Hinweis, dass die Kurven aus denselben Spielen und demselben
        Modell entstehen wie die Prognose, mit Verweis auf die Methodik. Er
        nennt bewusst keine Modellkonstanten.
  \item Die Tabelle \enquote{Stärkeschätzung heute}: Verein, ELO zum
        Saisonbeginn, ELO heute und Veränderung, sortiert nach der heutigen
        ELO. Sie wird auf dem Server erzeugt und funktioniert auch ohne
        JavaScript.
\end{enumerate}

\subsection*{Die Daten}

\code{elo\_verlauf\_daten()} (\datei{RCode/elo_verlauf.R:41--116}) baut für
jeden Verein eine Punktfolge. Punkt 0 ist die ELO zum Saisonstart. Danach
folgt je \emph{gespieltem} Spiel ein Punkt mit Datum, Spieltag,
Nachholspiel-Kennzeichen, Gegner, Ergebnis, ELO vor und nach dem Spiel, der
Änderung $\delta$ und den Wahrscheinlichkeiten für Sieg und Remis vor dem
Spiel. Gewertete Spiele erzeugen keinen Punkt, weil sie die ELO nicht
bewegen.

Eine Invariante sichert die Konsistenz mit der Ligatabelle: Der letzte Punkt
jeder Linie muss auf $10^{-6}$ genau der ELO in der Tabelle entsprechen.
Sonst entfällt diese eine Verlaufsseite mit einer Warnung im Log, und alle
anderen Seiten entstehen trotzdem.

Die Daten gehen als JSON in einem
\code{<script type="application/json">}-Element in die Seite. Jedes
\enquote{\textless} wird dabei als \code{\textbackslash u003c} kodiert. Das
bleibt gültiges JSON, und kein Vereinsname kann das Skript-Element mit
\code{</script>} oder \code{<!--} vorzeitig beenden.

\subsection*{Das Diagramm}

Das Skript \datei{RCode/site_assets/verlauf.js} zeichnet ein SVG mit dem
Koordinatensystem $1080 \times 640$.

\begin{description}[font=\sffamily\bfseries\small]
  \item[Wenig Tinte] Keine Rahmen, Achsenlinien, Schatten oder Flächen.
        Waagerechte Hilfslinien liegen alle 50 ELO-Punkte. Alle Linien sind
        leise grau, genau eine ist hervorgehoben (Tintenblau, 2~px). Es gibt
        keine Vereinsfarben und keine Legende. Jede Linie ist an ihrem Ende
        direkt beschriftet, mit dem Kürzel und der gerundeten heutigen ELO.
  \item[Senkrechte Skala] Sie reicht vom kleinsten bis zum größten Wert, mit
        einem Polster von 6~Prozent der Spanne, mindestens 10 Punkte. Liegt
        die Oberkante nah an einer 50er-Linie, rastet sie dort ein, damit das
        Diagramm oben mit einer Hilfslinie abschließt.
  \item[Zwei Zeitachsen] \enquote{Nach Spielen} (Standard) setzt das
        $n$-te Spiel jedes Vereins an dieselbe Stelle. \enquote{Nach Datum}
        zeigt die echte Zeitachse mit Monatsnamen und markiert die
        Winterpause als Fläche, wenn die größte Lücke länger als 18 Tage ist.
        Bei laufender Saison bleibt rechts Platz für den nächsten Spieltag.
  \item[Beschriftungen] Die Endbeschriftungen stehen in einer gemeinsamen
        Spalte rechts und sind auf mindestens 12{,}5~px Abstand
        auseinandergezogen. Vereine mit weniger Spielen erreichen die Spalte
        über eine gepunktete Führungslinie.
  \item[Tooltip] Zeigen oder Tippen auf eine Linie wählt den Verein und das
        letzte Spiel links vom Zeiger. Der Tooltip zeigt Datum und Spieltag,
        Paarung und Ergebnis, \enquote{Das Modell gab \emph{X} vorher
        $a$\,\% Sieg, $b$\,\% Remis.} und \enquote{ELO vor → nach $\pm\delta$}.
  \item[Tabelle und Diagramm] Fährt die Maus über eine Tabellenzeile, wird
        die zugehörige Linie hervorgehoben. Ein Klick hält die Auswahl.
        Umgekehrt markiert die Auswahl im Diagramm die Tabellenzeile.
  \item[Telefon] Das Diagramm bleibt mindestens 640~px breit und lässt sich
        seitlich scrollen (\code{touch-action: pan-x pan-y}).
\end{description}
